%!TEX TS-program = xelatex
%!TEX encoding = UTF-8 Unicode
% Awesome CV LaTeX Template for CV/Resume
%
% This template has been downloaded from:
% https://github.com/posquit0/Awesome-CV
%
% Author:
% Claud D. Park <posquit0.bj@gmail.com>
% http://www.posquit0.com
%
%
% Adapted to be an Rmarkdown template by Mitchell O'Hara-Wild
% 23 November 2018
%
% Template license:
% CC BY-SA 4.0 (https://creativecommons.org/licenses/by-sa/4.0/)
%
%-------------------------------------------------------------------------------
% CONFIGURATIONS
%-------------------------------------------------------------------------------
% A4 paper size by default, use 'letterpaper' for US letter
\documentclass[11pt,a4paper,]{awesome-cv}

% Configure page margins with geometry
\usepackage{geometry}
\geometry{left=1.4cm, top=.8cm, right=1.4cm, bottom=1.8cm, footskip=.5cm}


% Specify the location of the included fonts
\fontdir[fonts/]

% Color for highlights
% Awesome Colors: awesome-emerald, awesome-skyblue, awesome-red, awesome-pink, awesome-orange
%                 awesome-nephritis, awesome-concrete, awesome-darknight

\definecolor{awesome}{HTML}{414141}

% Colors for text
% Uncomment if you would like to specify your own color
% \definecolor{darktext}{HTML}{414141}
% \definecolor{text}{HTML}{333333}
% \definecolor{graytext}{HTML}{5D5D5D}
% \definecolor{lighttext}{HTML}{999999}

% Set false if you don't want to highlight section with awesome color
\setbool{acvSectionColorHighlight}{true}

% If you would like to change the social information separator from a pipe (|) to something else
\renewcommand{\acvHeaderSocialSep}{\quad\textbar\quad}

\def\endfirstpage{\newpage}

%-------------------------------------------------------------------------------
%	PERSONAL INFORMATION
%	Comment any of the lines below if they are not required
%-------------------------------------------------------------------------------
% Available options: circle|rectangle,edge/noedge,left/right

\name{Thiago}{Cerqueira Silva}

\position{Estudate-M.D./Ph.D.}
\address{Av. Reitor Miguel Calmon, S/N - Vale do Canela, Salvador-Bahia}

\mobile{+55 71-99124-7896}
\email{\href{mailto:csilvathiago@hotmail.com}{\nolinkurl{csilvathiago@hotmail.com}}}
\orcid{0000-0003-4534-2509}

% \gitlab{gitlab-id}
% \stackoverflow{SO-id}{SO-name}
% \skype{skype-id}
% \reddit{reddit-id}

\quote{Estudante de Medicina da Faculdade de Medicina da Bahia-UFBA e
doutorando do Programa de Pós Graduação em Ciências da Saúde (Programa
Pesquisador em Medicina). Experiência em trabalhos na área de
epidemiologia, bioestatística e bioinformática. Meus principais
interesses no campo acadêmico são métodos quantitativos, epidemiologia e
reprodutibilidade. Durante minha formação busquei obter conhecimentos
que me permitissem atuar na interface entre ciência de dados e Medicina,
utilizando os conhecimentos da prática clínica para um melhor
desenvolvimento de desenho de estudos, aplicação de métodos estatísticos
e interpretação de resultados. Durante a graduação participei na
elaboração e condução de diversos trabalhos científicos, incluindo
artigos na área de doença de Chagas, Chikungunya e COVID-19}

\usepackage{booktabs}

\providecommand{\tightlist}{%
	\setlength{\itemsep}{0pt}\setlength{\parskip}{0pt}}

%------------------------------------------------------------------------------


\usepackage{fancyhdr}
\pagestyle{fancy}
\fancyhead[R]{\thepage}

% Pandoc CSL macros
\newlength{\cslhangindent}
\setlength{\cslhangindent}{1.5em}
\newlength{\csllabelwidth}
\setlength{\csllabelwidth}{3em}
\newenvironment{CSLReferences}[3] % #1 hanging-ident, #2 entry spacing
 {% don't indent paragraphs
  \setlength{\parindent}{0pt}
  % turn on hanging indent if param 1 is 1
  \ifodd #1 \everypar{\setlength{\hangindent}{\cslhangindent}}\ignorespaces\fi
  % set entry spacing
  \ifnum #2 > 0
  \setlength{\parskip}{#2\baselineskip}
  \fi
 }%
 {}
\usepackage{calc}
\newcommand{\CSLBlock}[1]{#1\hfill\break}
\newcommand{\CSLLeftMargin}[1]{\parbox[t]{\csllabelwidth}{#1}}
\newcommand{\CSLRightInline}[1]{\parbox[t]{\linewidth - \csllabelwidth}{#1}}
\newcommand{\CSLIndent}[1]{\hspace{\cslhangindent}#1}

\begin{document}

% Print the header with above personal informations
% Give optional argument to change alignment(C: center, L: left, R: right)
\makecvheader

% Print the footer with 3 arguments(<left>, <center>, <right>)
% Leave any of these blank if they are not needed
% 2019-02-14 Chris Umphlett - add flexibility to the document name in footer, rather than have it be static Curriculum Vitae
\makecvfooter
  {January 20, 2022}
    {Thiago Cerqueira Silva~~~·~~~Curriculum Vitae}
  {\thepage}


%-------------------------------------------------------------------------------
%	CV/RESUME CONTENT
%	Each section is imported separately, open each file in turn to modify content
%------------------------------------------------------------------------------



\hypertarget{formauxe7uxe3o-acaduxeamica}{%
\section{\texorpdfstring{\faIcon{university} Formação
acadêmica}{ Formação acadêmica}}\label{formauxe7uxe3o-acaduxeamica}}

\begin{cventries}
    \cventry{Medicina}{Universidade Federal da Bahia}{Faculdade de Medicina da Bahia}{2014--2022}{\begin{cvitems}
\item Monografia: Influência da apneia obstrutiva do sono sobre índices de vasorreatividade cerebral em pacientes com acidente vascular cerebral.
\end{cvitems}}
    \cventry{Programa Pesquisador em Medicina}{Universidade Federal da Bahia}{Pós-Graduação em Ciências da Saúde}{2017--2022}{\begin{cvitems}
\item Tese: Fatores associados a progressão e a patogênese da doença Chikungunya na fase crônica
\end{cvitems}}
\end{cventries}

\hypertarget{artigos-publicados}{%
\section{\texorpdfstring{\faIcon{newspaper} Artigos
Publicados}{ Artigos Publicados}}\label{artigos-publicados}}

\begin{cventries}
    \cventry{Katikireddi, S. V., \textbf{Cerqueira-Silva, T.}, Vasileiou, E., Robertson, C., Amele, S., Pan, J., … Sheikh, A. }{Two-dose ChAdOx1 nCoV-19 vaccine protection against COVID-19 hospital admissions and deaths over time: a retrospective, population-based cohort study in Scotland and Brazil}{The Lancet}{2022}{}\vspace{-4.0mm}
    \cventry{\textbf{Cerqueira-Silva, T.}, Oliveira, V. de A., Boaventura, V. S., Pescarini, J. M., Júnior, J. B., Machado, T. M., … Barral-Netto, M. }{Influence of age on the effectiveness and duration of protection of Vaxzevria and CoronaVac vaccines: A population-based study}{The Lancet Regional Health - Americas}{2022}{}\vspace{-4.0mm}
    \cventry{Bonyek-Silva, I., Machado, A. F. A., \textbf{Cerqueira-Silva, T.}, Nunes, S., Silva Cruz, M. R., Silva, J., … Tavares, N. M. }{LTB4-Driven Inflammation and Increased Expression of ALOX5/ACE2 During Severe COVID-19 in Individuals With Diabetes}{Diabetes}{2021}{}\vspace{-4.0mm}
    \cventry{Garcia, K. O., \textbf{Silva, T. C.}, Nunes, M. C. P., Felix, I. F., Oliveira, M. A., Marques, M. E. L., … Oliveira-Filho, J. }{Tissue Inhibitor of Metalloproteinase-1 Is Increased in Chagasic Cardiomyopathy}{The American Journal of Tropical Medicine and Hygiene}{2021}{}\vspace{-4.0mm}
    \cventry{Silveira, M. A. D., De Jong, D., Berretta, A. A., Galvão, E. B. dos S., Ribeiro, J. C., \textbf{Cerqueira-Silva, T.}, … Passos, R. da H. }{Efficacy of Brazilian green propolis (EPP-AF®) as an adjunct treatment for hospitalized COVID-19 patients: A randomized, controlled clinical trial}{Biomedicine \& Pharmacotherapy}{2021}{}\vspace{-4.0mm}
    \cventry{\textbf{Cerqueira-Silva, T.}, Gonçalves, B. M., Pereira, C. B., Porto, L. M., Marques, M. E., Santos, L. S., … Oliveira-Filho, J. }{Chagas disease is an independent predictor of stroke and death in a cohort of heart failure patients}{International Journal of Stroke}{2021}{}\vspace{-4.0mm}
    \cventry{\textbf{Cerqueira-Silva, T.}, Carreiro, R., Nunes, V., Passos, L., Canedo, B. F., Andrade, S., … Boaventura, V. }{Bridging Learning in Medicine and Citizenship During the COVID-19 Pandemic: A Telehealth-Based Case Study}{JMIR Public Health and Surveillance}{2021}{}\vspace{-4.0mm}
    \cventry{Dias, A. R., Fernandes, S. M., Fialho-Silva, I., \textbf{Cerqueira-Silva, T.}, Miranda-Scippa, A., \& Galvão-de Almeida, A. }{Burnout syndrome and resilience in medical students from a Brazilian public college in Salvador, Brazil}{Trends in Psychiatry and Psychotherapy}{2021}{}\vspace{-4.0mm}
    \cventry{De Moraes, L., \textbf{Cerqueira-Silva, T.}, Nobrega, V., Akrami, K., Santos, L. A., Orge, C., … Boaventura, V. S. }{A clinical scoring system to predict long-term arthralgia in Chikungunya disease: A cohort study}{PLOS Neglected Tropical Diseases}{2020}{}\vspace{-4.0mm}
    \cventry{Casais, P. M., Akrami, K., \textbf{Cerqueira-Silva, T.}, Moraes, L. P., Rigaud, V. N., Neto, E. S., … Boaventura, V. S. }{Oral lesions are frequent in patients with Chikungunya infection}{Journal of Travel Medicine}{2020}{}\vspace{-4.0mm}
    \cventry{Castello‐Branco, R. C., Cerqueira‐Silva, T., Andrade, A. L., Gonçalves, B. M. M., Pereira, C. B., Felix, I. F., … Oliveira‐Filho, J. }{Association Between Risk of Obstructive Sleep Apnea and Cerebrovascular Reactivity in Stroke Patients}{Journal of the American Heart Association}{2020}{}\vspace{-4.0mm}
    \cventry{Neri, L. R., Torreão, J. A., Porto, L. M., Gonçalves, B. M., Andrade, A. L., Pereira, C. B., … Oliveira-Filho, J. }{Factors associated with abnormal cardiac magnetic resonance imaging in embolic stroke of undetermined source}{International Journal of Stroke}{2019}{}\vspace{-4.0mm}
\end{cventries}

\hypertarget{artigos-aceitos-para-publicauxe7uxe3o}{%
\section{\texorpdfstring{\faIcon{newspaper} Artigos aceitos para
publicação}{ Artigos aceitos para publicação}}\label{artigos-aceitos-para-publicauxe7uxe3o}}

\begin{cventries}
    \cventry{\textbf{Cerqueira-Silva, T.} Katikireddi, S. V., ...Barral-Netto, Manoel; Boaventura, Viviane Sampaio}{Vaccine Effectiveness of heterologous CoronaVac plus BNT162b2 in Brazil}{Nature Medicine}{2022}{}\vspace{-4.0mm}
\end{cventries}

\hypertarget{preprints}{%
\section{\texorpdfstring{\faIcon{newspaper}
Preprints}{ Preprints}}\label{preprints}}

\begin{cventries}
    \cventry{\textbf{Cerqueira-Silva, T.}, Andrews, J. R., Boaventura, V. S., Ranzani, O. T., de Araújo Oliveira, V., Paixão, E. S., … Barral-Netto, M. }{Effectiveness of CoronaVac, ChAdOx1, BNT162b2 and Ad26.COV2.S among individuals with prior SARS-CoV-2 infection in Brazil}{medRxiv}{2021}{}\vspace{-4.0mm}
    \cventry{Andrade Barreto, A. P., Duarte, L. C., \textbf{Cerqueira-Silva, T.}, Barreto Filho, M. A., Camelier, A., Tavares, N. M., … Coelho Lima, M. C. }{Post-Acute COVID Syndrome, the Aftermath of Mild to Severe COVID-19 in Brazilian Patients}{medRxiv}{2021}{}\vspace{-4.0mm}
    \cventry{Paixão, E. S., Wong, K. L., Alves, F. J. O., de Araújo Oliveira, V., \textbf{Cerqueira-Silva, T.}, Júnior, J. B., … Barral-Netto, M. }{Effectiveness of the CoronaVac Vaccine in Prevention of Symptomatic and Progression to Severe COVID-19 in Pregnant Women in Brazil}{SSRN Electronic Journal}{2021}{}\vspace{-4.0mm}
\end{cventries}

\hypertarget{experiuxeancia-em-pesquisa}{%
\section{\texorpdfstring{\faIcon{microscope} Experiência em
Pesquisa}{ Experiência em Pesquisa}}\label{experiuxeancia-em-pesquisa}}

\begin{cventries}
    \cventry{Atuo no projeto que estuda a efetividade das vacinas contra COVID-19.}{VigiVac}{Fundação Oswaldo Cruz}{2021-Atual}{}\vspace{-4.0mm}
    \cventry{Atuo no projeto que estuda os fatores envolvidos na persistência de artralgia em pacientes pós-doença Chikungunya, tanto nas características clínicas como assinatura molecular}{LEITV}{Instituto Gonçalo Moniz-Fiocruz Bahia}{2017-Atual}{}\vspace{-4.0mm}
    \cventry{Atuo no projeto que estuda a associação da doença de Chagas e acidente vascular isquêmico em pacientes com insuficiência cardíaca.}{Grupo Neurocardio}{FAMEB-UFBA}{2016-Atual}{}\vspace{-4.0mm}
\end{cventries}

\newpage

\hypertarget{bolsa-de-estudos}{%
\section{\texorpdfstring{\faIcon{money-check-alt} Bolsa de
Estudos}{ Bolsa de Estudos}}\label{bolsa-de-estudos}}

\begin{cventries}
    \cventry{}{O uso de aprendizado de máquina para diagnóstico da COVID-19}{UFBA}{2021--2022}{}\vspace{-4.0mm}
    \cventry{}{Monitoramento e avaliação das ações estratégicas dos programas de provimento e formação em saúde do ministério da saúde.}{Fiotec}{2020--2021}{}\vspace{-4.0mm}
    \cventry{}{Expressão gênica nas lesões de leishmaniose cutânea e suas variantes clínicas através de técnicas de alto rendimento}{Fiotec}{2019--2020}{}\vspace{-4.0mm}
    \cventry{}{CHADSS: Chagas Disease Scan Study Neuroimagem na Doença de Chagas}{NIH}{2016--2017}{}\vspace{-4.0mm}
\end{cventries}

\hypertarget{cursos}{%
\section{\texorpdfstring{\faIcon{laptop-code}
Cursos}{ Cursos}}\label{cursos}}

\begin{cventries}
    \cventry{Federal University of Minas Gerais}{III Curso de Verão em Bioinformática da UFMG}{}{2019}{}\vspace{-4.0mm}
    \cventry{Fundação Oswaldo Cruz - Instituto Carlos Chagas}{II International Course on Bioinformatics Analysis of RNA-Seq Data}{}{2019}{}\vspace{-4.0mm}
    \cventry{Instituto Gonçalo Moniz - Fiocruz Bahia}{XVII Curso Internacional de Epidemiologia Molecular em Doenças Infecciosas}{}{2017}{}\vspace{-4.0mm}
    \cventry{StanfordOnline}{STATSX0001: Statistical Learning}{}{2021}{}\vspace{-4.0mm}
    \cventry{Universidade Johns Hopkins-Coursera}{Reproducible Research}{}{2020}{}\vspace{-4.0mm}
    \cventry{Imperial College London-Coursera}{Statistical Analysis with R for Public Health - Specialization}{}{2020}{}\vspace{-4.0mm}
    \cventry{HarvardX}{PH125.8x: Data Science: Machine Learning}{}{2020}{}\vspace{-4.0mm}
    \cventry{Imperial College London-Coursera}{Epidemiology for Public Health  - Specialization}{}{2020}{}\vspace{-4.0mm}
    \cventry{Moscow Institute of Physics and Technology-Coursera}{Academic Literacy}{}{2019}{}\vspace{-4.0mm}
    \cventry{Universidade Johns Hopkins-Coursera}{R Programming}{}{2018}{}\vspace{-4.0mm}
    \cventry{HarvardX}{PH525.1x: Statistics and R}{}{2018}{}\vspace{-4.0mm}
\end{cventries}

\newpage

\hypertarget{monitoria}{%
\section{\texorpdfstring{\faIcon{user-edit}
Monitoria}{ Monitoria}}\label{monitoria}}

\begin{cventries}
    \cventry{Atuei como monitor da disciplina de Imunopatologia, auxiliando os alunos da disciplina na preparação de seminários e resolução de exercícios e contribuindo no processo avaliativo da matéria.}{MEB21-Imunopatologia I}{FMB-UFBA}{2019.1/2019.2}{}\vspace{-4.0mm}
\end{cventries}

\hypertarget{capuxedtulos-de-livros}{%
\section{\texorpdfstring{\faIcon{book} Capítulos de
livros}{ Capítulos de livros}}\label{capuxedtulos-de-livros}}

\begin{cventries}
    \cventry{Construção de conhecimento no curso da pandemia de COVID-19: aspectos biomédicos, clínico-assistenciais, epidemiológicos e sociais}{Vacinas contra a infecção pelo SARS-CoV-2}{}{2020}{}\vspace{-4.0mm}
    \cventry{Construção de conhecimento no curso da pandemia de COVID-19: aspectos biomédicos, clínico-assistenciais, epidemiológicos e sociais}{Testes diagnósticos na COVID-19}{}{2020}{}\vspace{-4.0mm}
    \cventry{Construção de conhecimento no curso da pandemia de COVID-19: aspectos biomédicos, clínico-assistenciais, epidemiológicos e sociais}{Aprendizado em medicina e cidadania durante a pandemia de COVID-19: um modelo de teletriagem em saúde}{}{2020}{}\vspace{-4.0mm}
\end{cventries}

\hypertarget{apresentauxe7uxf5es-em-congressos}{%
\section{\texorpdfstring{\faIcon{comments} Apresentações em
congressos}{ Apresentações em congressos}}\label{apresentauxe7uxf5es-em-congressos}}

\begin{cventries}
    \cventry{XII Congresso Brasileiro de Doenças Cerebrovasculares}{Coorte de 13 anos de pacientes com insuficiência cardíaca chagásica e não-chagásica: fatores de risco para AVC e morte}{}{2019}{}\vspace{-4.0mm}
\end{cventries}

\hypertarget{resumos-publicados-em-anais-de-congressos}{%
\section{\texorpdfstring{\faIcon{comments} Resumos publicados em anais
de
congressos}{ Resumos publicados em anais de congressos}}\label{resumos-publicados-em-anais-de-congressos}}

\begin{cventries}
    \cventry{29th European Congress of Psychiatry}{Lithium treatment and estimate glomerular filtration
rate in bipolar disorder patients: A cross-sectional study}{}{2021}{}\vspace{-4.0mm}
    \cventry{5th European Stroke Organisation Conference}{Influence of risk of obstructive sleep apnea on cerebral vasoreactivity in patients with stroke}{}{2019}{}\vspace{-4.0mm}
    \cventry{XII Congresso Brasileiro de Doenças Cerebrovasculares}{Coorte de 13 anos de pacientes com insuficiência cardíaca chagásica e não-chagásica: fatores de risco para AVC e morte}{}{2019}{}\vspace{-4.0mm}
    \cventry{XII Congresso Brasileiro de Doenças Cerebrovasculares}{Mecanismos de disfunção cognitiva em pacientes com insuficiência cardíaca chagásica e não-chagásica}{}{2019}{}\vspace{-4.0mm}
    \cventry{XII Congresso Brasileiro de Doenças Cerebrovasculares}{Depressão e ansiedade em portadores de insuficiência cardíaca de etiologia chagásica e não- chagásica. In: Depressão e ansiedade em portadores de insuficiência cardíaca de etiologia chagásica e não- chagásica}{}{2019}{}\vspace{-4.0mm}
    \cventry{XII Congresso Brasileiro de Doenças Cerebrovasculares}{Etiologias de acidente vascular cerebral isquêmico em pacientes com doença de chagas}{}{2019}{}\vspace{-4.0mm}
    \cventry{11th World Stroke Congress}{Neuronal depopulation, not silent cerebrovascular disease mediate cognitive dysfunction in Chagas disease}{}{2018}{}\vspace{-4.0mm}
    \cventry{XI Congresso Brasileiro de Doenças Cerebrovasculares}{Impacto da estenose intracraniana na taxa de AVC e morte em pacientes com acidente vascular cerebral isquêmico}{}{2017}{}\vspace{-4.0mm}
    \cventry{XI Congresso Brasileiro de Doenças Cerebrovasculares}{Preditores de acidente vascular cerebral e óbito em uma coorte de pacientes com cardiopatia chagásica e não chagásica}{}{2017}{}\vspace{-4.0mm}
    \cventry{XI Congresso Brasileiro de Doenças Cerebrovasculares}{Doença cerebrovascular silenciosa: preditores de disfunção cognitiva em portadores de cardiomiopatia chagásica. In: XI Congresso Brasileiro de Doenças Cerebrovasculares e XV Jornada de Neurologia da Bahia}{}{2017}{}\vspace{-4.0mm}
    \cventry{3rd European Stroke Organisation Conference}{Pilot study of cardiac magnetic resonance imaging in embolic stroke of undetermined source (MR-ESUS)}{}{2017}{}\vspace{-4.0mm}
\end{cventries}

\hypertarget{extensuxe3o-universituxe1ria}{%
\section{\texorpdfstring{\faIcon{hands-helping} Extensão
Universitária}{ Extensão Universitária}}\label{extensuxe3o-universituxe1ria}}

\begin{cventries}
    \cventry{Participei do projeto ``O uso de aprendizado de máquina para diagnóstico da COVID-19'' em que foi desenvolvido uma calculadora para predizer o diagnóstico de COVID-19 com base na sintomatologia do paciente.}{PIBIEX-UFBA}{}{Agosto-2021/Janeiro-2022}{}\vspace{-4.0mm}
    \cventry{Criação de textos de divulgação científica durante o início da pandemia de COVID-19 (2020)}{Rede CoVida}{}{Março-2020/Setembro-2020}{}\vspace{-4.0mm}
    \cventry{Participei da criação e coordenação do projeto Telecoronavírus, durante o início da pandemia da COVID-19 na Bahia.}{Telecoronavirus: Ação das Escolas e Entidades Médicas da Bahia.}{}{Março-2020/Julho-2020}{}\vspace{-4.0mm}
    \cventry{Participei da criação da liga acadêmica em ciências da Faculdade de Medicina em 2018.1}{Liga Acadêmica de Ciências}{}{Janeiro-2018/Janeiro-2022}{}\vspace{-4.0mm}
\end{cventries}

\hypertarget{notas}{%
\section{Notas}\label{notas}}

\begin{itemize}
\tightlist
\item
  Para uma versão mais detalhada do meu curriculum vitae, por favor
  visite meu \href{http://lattes.cnpq.br/8074214277248039}{Lattes}
\item
  Este currículo é totalmente reprodutível; todo o código utilizado para
  a produção dele está disponível neste
  \href{https://github.com/csthiago/cv}{repositório do GitHub}.
\end{itemize}

\includegraphics{sig/Assinatura_2.png}\\



\end{document}
